\begin{figure*}[t]
\centering
\setcounter{subfigure}{0}
\begin{minipage}{\textwidth}
\centering
\ESPanel{.24\linewidth}{panels/fig04/fig04_a.pdf}{Boundary shift: 0.4 m}%
\hfill%
\ESPanel{.24\linewidth}{panels/fig04/fig04_b.pdf}{Boundary shift: 1.2 m}%
\hfill%
\ESPanel{.24\linewidth}{panels/fig04/fig04_c.pdf}{Local height shift: $\pm 0.3$ m}%
\hfill%
\ESPanel{.24\linewidth}{panels/fig04/fig04_d.pdf}{Local height shift: $\pm 1.2$ m}

\par\vspace{2pt}
\includegraphics[width=.82\linewidth]{\FigureAssetRoot/panels/fig04/fig04_legend.pdf}
\par\vspace{2pt}

\end{minipage}
\caption{Controlled examples show how geometric changes affect Q. The reference room has a $4\times4$ m footprint, a 3 m wall height and a 1.6 m camera height. (a--b) Moving one floor boundary inward reduces overlap and increases boundary distance; hatched regions mark omitted floor area. (c--d) Opposing upper-boundary shifts preserve the footprint and perimeter-mean height while increasing local height error. Blue and orange show the smaller and larger changes; dashed gray denotes the reference. Measurements and Q are computed from the displayed geometry.}
\label{fig:quality-components}
\end{figure*}
